\documentclass[10pt,a4paper]{amsart}
\usepackage{amsmath,amssymb,amsthm,graphicx}
% Retained minimal preamble for the Stacks Project chapter flat.tex (More on Flatness).
% Verbatim excerpt of the chapter preamble: Xy-pic setup (native preamble lines 14-18) and the native
% theorem declarations with the mathematical macros (native preamble lines 154-217). Copied mechanically;
% nothing added, removed or redefined. The wrapper class and the amsmath/amssymb/amsthm/graphicx base
% packages are supplied by the assembler, so the class/externaldocument part of the native preamble is omitted.
% For commutative diagrams we use Xy-pic
\usepackage[all]{xy}

% We use 2cell for 2-commutative diagrams.
\xyoption{2cell}
% Theorem environments.
%
\theoremstyle{plain}
\newtheorem{theorem}[subsection]{Theorem}
\newtheorem{proposition}[subsection]{Proposition}
\newtheorem{lemma}[subsection]{Lemma}

\theoremstyle{definition}
\newtheorem{definition}[subsection]{Definition}
\newtheorem{example}[subsection]{Example}
\newtheorem{exercise}[subsection]{Exercise}
\newtheorem{situation}[subsection]{Situation}

\theoremstyle{remark}
\newtheorem{remark}[subsection]{Remark}
\newtheorem{remarks}[subsection]{Remarks}

\numberwithin{equation}{subsection}

% Macros
%
\def\lim{\mathop{\mathrm{lim}}\nolimits}
\def\colim{\mathop{\mathrm{colim}}\nolimits}
\def\Spec{\mathop{\mathrm{Spec}}}
\def\Hom{\mathop{\mathrm{Hom}}\nolimits}
\def\Ext{\mathop{\mathrm{Ext}}\nolimits}
\def\SheafHom{\mathop{\mathcal{H}\!\mathit{om}}\nolimits}
\def\SheafExt{\mathop{\mathcal{E}\!\mathit{xt}}\nolimits}
\def\Sch{\mathit{Sch}}
\def\Mor{\mathop{\mathrm{Mor}}\nolimits}
\def\Ob{\mathop{\mathrm{Ob}}\nolimits}
\def\Sh{\mathop{\mathit{Sh}}\nolimits}
\def\NL{\mathop{N\!L}\nolimits}
\def\CH{\mathop{\mathrm{CH}}\nolimits}
\def\proetale{{pro\text{-}\acute{e}tale}}
\def\etale{{\acute{e}tale}}
\def\QCoh{\mathit{QCoh}}
\def\Ker{\mathop{\mathrm{Ker}}}
\def\Im{\mathop{\mathrm{Im}}}
\def\Coker{\mathop{\mathrm{Coker}}}
\def\Coim{\mathop{\mathrm{Coim}}}

% Boxtimes
%
\DeclareMathSymbol{\boxtimes}{\mathbin}{AMSa}{"02}

%
% Macros for moduli stacks/spaces
%
\def\QCohstack{\mathcal{QC}\!\mathit{oh}}
\def\Cohstack{\mathcal{C}\!\mathit{oh}}
\def\Spacesstack{\mathcal{S}\!\mathit{paces}}
\def\Quotfunctor{\mathrm{Quot}}
\def\Hilbfunctor{\mathrm{Hilb}}
\def\Curvesstack{\mathcal{C}\!\mathit{urves}}
\def\Polarizedstack{\mathcal{P}\!\mathit{olarized}}
\def\Complexesstack{\mathcal{C}\!\mathit{omplexes}}
% \Pic is the operator that assigns to X its picard group, usage \Pic(X)
% \Picardstack_{X/B} denotes the Picard stack of X over B
% \Picardfunctor_{X/B} denotes the Picard functor of X over B
\def\Pic{\mathop{\mathrm{Pic}}\nolimits}
\def\Picardstack{\mathcal{P}\!\mathit{ic}}
\def\Picardfunctor{\mathrm{Pic}}
\def\Deformationcategory{\mathcal{D}\!\mathit{ef}}

% --- appended 2026-10-12 (Chief): the project own definitions of the reference/slogan/history
% display environments, copied mechanically from the shared Stacks preamble lines 8-12
% (\usepackage{verbatim} plus the three \newenvironment lines, which the chapter preamble keeps
% verbatim from the project preamble). Nothing added, removed or redefined: these environments
% carry no mathematical content and the project itself typesets them as comments. Without these
% three lines a counted statement whose own text uses them fails with Environment slogan undefined.
\usepackage{verbatim}
\newenvironment{reference}{\comment}{\endcomment}
\newenvironment{slogan}{\comment}{\endcomment}
\newenvironment{history}{\comment}{\endcomment}
\usepackage{cleveref}
\crefname{originaldefinition}{original definition}{original definitions}
\Crefname{originaldefinition}{Original definition}{Original definitions}
\makeatletter
\newlabel{algebra-definition-henselization}{{04GQ}{source}{}{sourcepointer-1}{}}
\newlabel{algebra-definition-henselization@cref}{{[originaldefinition][0][]04GQ}{[1][0][]source}{}{}{}}
\makeatother
\crefname{originallemma}{original lemma}{original lemmas}
\Crefname{originallemma}{Original lemma}{Original lemmas}
\makeatletter
\newlabel{algebra-lemma-faithfully-flat-universally-injective}{{05CK}{source}{}{sourcepointer-2}{}}
\newlabel{algebra-lemma-faithfully-flat-universally-injective@cref}{{[originallemma][0][]05CK}{[1][0][]source}{}{}{}}
\makeatother
\crefname{originallemma}{original lemma}{original lemmas}
\Crefname{originallemma}{Original lemma}{Original lemmas}
\makeatletter
\newlabel{algebra-lemma-strict-henselization-different}{{04GW}{source}{}{sourcepointer-3}{}}
\newlabel{algebra-lemma-strict-henselization-different@cref}{{[originallemma][0][]04GW}{[1][0][]source}{}{}{}}
\makeatother
\crefname{originallemma}{original lemma}{original lemmas}
\Crefname{originallemma}{Original lemma}{Original lemmas}
\makeatletter
\newlabel{algebra-lemma-weakly-ass-local}{{0566}{source}{}{sourcepointer-4}{}}
\newlabel{algebra-lemma-weakly-ass-local@cref}{{[originallemma][0][]0566}{[1][0][]source}{}{}{}}
\makeatother
\crefname{originaldefinition}{original definition}{original definitions}
\Crefname{originaldefinition}{Original definition}{Original definitions}
\makeatletter
\newlabel{divisors-definition-weakly-associated}{{056L}{source}{}{sourcepointer-5}{}}
\newlabel{divisors-definition-weakly-associated@cref}{{[originaldefinition][0][]056L}{[1][0][]source}{}{}{}}
\makeatother
\crefname{originallemma}{original lemma}{original lemmas}
\Crefname{originallemma}{Original lemma}{Original lemmas}
\makeatletter
\newlabel{divisors-lemma-weakly-ass-pullback}{{05F0}{source}{}{sourcepointer-6}{}}
\newlabel{divisors-lemma-weakly-ass-pullback@cref}{{[originallemma][0][]05F0}{[1][0][]source}{}{}{}}
\makeatother
\crefname{originallemma}{original lemma}{original lemmas}
\Crefname{originallemma}{Original lemma}{Original lemmas}
\makeatletter
\newlabel{lemma-etale-weak-assassin-up-down}{{05FP}{source}{}{sourcepointer-7}{}}
\newlabel{lemma-etale-weak-assassin-up-down@cref}{{[originallemma][0][]05FP}{[1][0][]source}{}{}{}}
\makeatother
\crefname{originallemma}{original lemma}{original lemmas}
\Crefname{originallemma}{Original lemma}{Original lemmas}
\makeatletter
\newlabel{more-algebra-lemma-dumb-properties-henselization}{{07QM}{source}{}{sourcepointer-8}{}}
\newlabel{more-algebra-lemma-dumb-properties-henselization@cref}{{[originallemma][0][]07QM}{[1][0][]source}{}{}{}}
\makeatother
\title{Weakly associated points are unchanged under henselization and strict henselization}
\author{The Stacks Project Authors}
\date{stacks-project:flat.tex:local-snapshot-20261011; GNU FDL 1.2 or later}
\begin{document}
\maketitle
\noindent\textit{Editorial note.} All delivered mathematical text is the author's own, copied from the pinned Stacks Project chapter file flat.tex; no statement or proof step was rewritten, reordered or removed.\par\medskip
\noindent\textit{Editorial note.} Preamble provenance: the chapter's own retained minimal preamble (sources/stacks-flat/retained-preamble.tex, lines 1-85) is a verbatim excerpt of the native Xy-pic setup and of the native theorem declarations and mathematical macros, followed by the project's own reference/slogan/history display-environment definitions copied mechanically from the shared preamble; the wrapper class and the amsmath/amssymb/amsthm/graphicx base packages are supplied by the assembler, and extra\_packages is left empty because everything needed is already in the retained preamble. No macro was invented or redefined.\par\medskip
\noindent\textit{Editorial note.} Same-chapter pointers are located at their \textbackslash{}label lines inside the pinned chapter file; cross-chapter pointers are located in the project's own tag list staged in the snapshot (sources/stacks-chow/tags.txt), at the verbatim lines carrying their canonical tags.\par\medskip
\noindent\textit{Allowed standard prerequisites (editorial).}\begin{itemize}
\item The definition of a weakly associated point of a module. Divisors, label divisors-definition-weakly-associated, invoked at proof line 399. Gives the equivalence of the first two conditions, which is just the definition rewritten at the level of the local ring.
\item Behaviour of weakly associated primes under a flat morphism, and faithful flatness of the henselization maps. Divisors, label divisors-lemma-weakly-ass-pullback (proof line 401); More on Algebra, label more-algebra-lemma-dumb-properties-henselization (proof line 403). The first gives the implications from the local ring to the henselization and then to the strict henselization, since the two localisation maps are flat, and the second identifies those maps as faithfully flat.
\item The local description of weakly associated primes, the colimit description of the strict henselization, and the up and down lemma for an etale neighbourhood. Algebra, labels algebra-lemma-weakly-ass-local (proof line 414), algebra-lemma-strict-henselization-different (proof line 418), algebra-lemma-faithfully-flat-universally-injective (proof line 424); Same chapter flat.tex, label lemma-etale-weak-assassin-up-down (proof line 429). Turns the remaining implication into a statement about an element whose annihilator has radical the maximal ideal, descends that element to an etale neighbourhood through the colimit description, and transfers the conclusion back along the etale morphism by the up and down lemma.
\end{itemize}\medskip
\ifdefined\hypertarget\hypertarget{sourcepointer-1}{}\fi
\noindent\textit{Source reference.} Original definition 04GQ: sources/stacks-chow/tags.txt, source lines 5841--5841. Navigation pointer for \textbackslash{}ref\{algebra-definition-henselization\} used inside the counted span (line 384); the argument that consumes it is retained in the span.\par\medskip
\ifdefined\hypertarget\hypertarget{sourcepointer-2}{}\fi
\noindent\textit{Source reference.} Original lemma 05CK: sources/stacks-chow/tags.txt, source lines 6923--6923. Navigation pointer for \textbackslash{}ref\{algebra-lemma-faithfully-flat-universally-injective\} used inside the counted span (line 424); the argument that consumes it is retained in the span.\par\medskip
\ifdefined\hypertarget\hypertarget{sourcepointer-3}{}\fi
\noindent\textit{Source reference.} Original lemma 04GW: sources/stacks-chow/tags.txt, source lines 5847--5847. Navigation pointer for \textbackslash{}ref\{algebra-lemma-strict-henselization-different\} used inside the counted span (line 418); the argument that consumes it is retained in the span.\par\medskip
\ifdefined\hypertarget\hypertarget{sourcepointer-4}{}\fi
\noindent\textit{Source reference.} Original lemma 0566: sources/stacks-chow/tags.txt, source lines 6699--6699. Navigation pointer for \textbackslash{}ref\{algebra-lemma-weakly-ass-local\} used inside the counted span (line 414); the argument that consumes it is retained in the span.\par\medskip
\ifdefined\hypertarget\hypertarget{sourcepointer-5}{}\fi
\noindent\textit{Source reference.} Original definition 056L: sources/stacks-chow/tags.txt, source lines 6714--6714. Navigation pointer for \textbackslash{}ref\{divisors-definition-weakly-associated\} used inside the counted span (line 399); the argument that consumes it is retained in the span.\par\medskip
\ifdefined\hypertarget\hypertarget{sourcepointer-6}{}\fi
\noindent\textit{Source reference.} Original lemma 05F0: sources/stacks-chow/tags.txt, source lines 7009--7009. Navigation pointer for \textbackslash{}ref\{divisors-lemma-weakly-ass-pullback\} used inside the counted span (line 401); the argument that consumes it is retained in the span.\par\medskip
\ifdefined\hypertarget\hypertarget{sourcepointer-7}{}\fi
\noindent\textit{Source reference.} Original lemma 05FP: sources/stacks-flat/pinned-original-source.tex, source lines 353--353. Navigation pointer for \textbackslash{}ref\{lemma-etale-weak-assassin-up-down\} used inside the counted span (line 429); the argument that consumes it is retained in the span.\par\medskip
\ifdefined\hypertarget\hypertarget{sourcepointer-8}{}\fi
\noindent\textit{Source reference.} Original lemma 07QM: sources/stacks-chow/tags.txt, source lines 9837--9837. Navigation pointer for \textbackslash{}ref\{more-algebra-lemma-dumb-properties-henselization\} used inside the counted span (line 403); the argument that consumes it is retained in the span.\par\medskip
\begin{lemma}
\label{lemma-weakly-associated-henselization}
Let $S$ be a scheme and $s \in S$ a point. Denote $\mathcal{O}_{S, s}^h$
(resp.\ $\mathcal{O}_{S, s}^{sh}$) the henselization (resp.\ strict
henselization), see
Algebra, Definition \ref{algebra-definition-henselization}.
Let $\mathcal{F}$ be a quasi-coherent $\mathcal{O}_S$-module.
The following are equivalent
\begin{enumerate}
\item $s$ is a weakly associated point of $\mathcal{F}$,
\item $\mathfrak m_s$ is a weakly associated prime of $\mathcal{F}_s$,
\item $\mathfrak m_s^h$ is a weakly associated prime of
$\mathcal{F}_s \otimes_{\mathcal{O}_{S, s}} \mathcal{O}_{S, s}^h$, and
\item $\mathfrak m_s^{sh}$ is a weakly associated prime of
$\mathcal{F}_s \otimes_{\mathcal{O}_{S, s}} \mathcal{O}_{S, s}^{sh}$.
\end{enumerate}
\end{lemma}

\begin{proof}
The equivalence of (1) and (2) is the definition, see
Divisors, Definition \ref{divisors-definition-weakly-associated}.
The implications (2) $\Rightarrow$ (3) $\Rightarrow$ (4)
follows from Divisors, Lemma \ref{divisors-lemma-weakly-ass-pullback}
applied to the flat (More on Algebra, Lemma
\ref{more-algebra-lemma-dumb-properties-henselization})
morphisms
$$
\Spec(\mathcal{O}_{S, s}^{sh}) \to
\Spec(\mathcal{O}_{S, s}^h) \to
\Spec(\mathcal{O}_{S, s})
$$
and the closed points. To prove (4) $\Rightarrow$ (2) we may replace
$S$ by an affine neighbourhood. Suppose that
$x \in \mathcal{F}_s \otimes_{\mathcal{O}_{S, s}} \mathcal{O}_{S, s}^{sh}$
is an element whose annihilator has radical equal to $\mathfrak m_s^{sh}$.
(See Algebra, Lemma \ref{algebra-lemma-weakly-ass-local}.)
Since $\mathcal{O}_{S, s}^{sh}$ is equal to the limit
of $\mathcal{O}_{U, u}$ over \'etale neighbourhoods
$f : (U, u) \to (S, s)$ by Algebra, Lemma
\ref{algebra-lemma-strict-henselization-different}
we may assume that $x$ is the image of some
$x' \in \mathcal{F}_s \otimes_{\mathcal{O}_{S, s}} \mathcal{O}_{U, u}$.
The local ring map $\mathcal{O}_{U, u} \to \mathcal{O}_{S, s}^{sh}$
is faithfully flat (as it is the strict henselization), hence
universally injective
(Algebra, Lemma \ref{algebra-lemma-faithfully-flat-universally-injective}).
It follows that the annihilator of $x'$ is the inverse image of the
annihilator of $x$. Hence the radical of this annihilator
is equal to $\mathfrak m_u$.
Thus $u$ is a weakly associated point of $f^*\mathcal{F}$.
By Lemma \ref{lemma-etale-weak-assassin-up-down}
we see that $s$ is a weakly associated point of $\mathcal{F}$.
\end{proof}

\end{document}
